\section{Discussion}

The paper's story is now sharper than the original draft. Under deployment shift, worst-group auditing is necessary because aggregate calibration can hide real subgroup failures. But under covariate shift, worst-group auditing is also statistically fragile because the relevant information is local. A subgroup-bin cell can become effectively unsupported long before the weighted sample looks globally unusable.

This perspective clarifies three points that are often conflated.

\begin{itemize}
  \item \textbf{Transport is not the same as stability.} Covariate shift tells us how to rewrite target quantities as weighted source quantities. It does not tell us whether those quantities can be estimated reliably in every audited cell.
  \item \textbf{Global ESS is not a worst-group diagnostic.} It summarizes the full weighted sample. Worst-group auditing depends on the weakest audited cell, so the relevant object is the lower tail of local ESS.
  \item \textbf{A subgroup alarm is not automatically evidence of a real deployment failure.} It may reflect true subgroup harm, or it may reflect local overlap collapse. The audit result is only interpretable once both possibilities are on the table.
\end{itemize}

This leads to a practical recommendation that is easy to state:

\begin{quote}
report worst-group reliability together with local overlap diagnostics.
\end{quote}

In concrete terms, a subgroup audit under covariate shift should not stop at one global ESS value. It should also report cell-level or lower-tail local ESS summaries, ideally for the exact subgroup-bin family used by the audit. Otherwise a reviewer or practitioner cannot tell whether a large worst-group estimate reflects deployment harm or statistical degeneracy.

There are several directions that could make the paper even stronger.

\begin{enumerate}
  \item Replace the finite subgroup family by a richer class and derive complexity-sensitive local-overlap bounds.
  \item Move from point estimates to interval estimates, for example confidence intervals for subgroup-bin gaps or worst-group error conditioned on local ESS.
  \item Study adaptive auditing, where the subgroup family itself is chosen after looking at the data, which should make false-alarm control even harder.
  \item Extend the local-overlap analysis beyond pure covariate shift to mixed shift, concept drift, and settings with doubly robust correction \citep{Yang2024DRCalibration, Wu2024MCOOD}.
\end{enumerate}

The broader claim is not about one library or one benchmark suite. It is a statistical claim about deployment auditing. If worst-group reliability is going to be used as evidence in shifted environments, then local overlap has to be part of the evidence. That is the main message of the paper.
